\scanpage{70}
\mspage{- 7 -}\hfill n\textsuperscript{o} \handwritten{380}

Dans le premier cas, on observe que $\lVert e-x_n^{-1}x\rVert=\lVert x_n^{-1}(x_n-x)\rVert\handwritten{\le}\lVert x_n^{-1}\rVert\ \lVert x_n-x\rVert$ donc tend vers zéro lorsque $n\handwritten{\to}\handwritten{\infty}$. Par suite, pour $n$ assez grand, $x_n^{-1}x$ est inversible, ce qui implique que $x$ est inversible.

Dans le second cas, compte tenu de la prop.2, e), on a $\handwritten{\unsure{\mathscr{D}|\mathfrak{D}|\partial}}(e-x_n^{-1}x)\handwritten{\le}\handwritten{\mathscr{D}}(x_n^{-1})\,\handwritten{\mathscr{D}}(x_n-x)\handwritten{\le}\handwritten{\mathscr{D}}(x_n^{-1})\ \lVert x_n-x\rVert$, donc $\handwritten{\mathscr{D}}(e-x_n^{-1}x)$ tend vers 0 lorsque $n\handwritten{\to}\handwritten{\infty}$, ce qui implique que pour $n$ assez grand $x_n^{-1}x$ est inversible et par conséquent, $x$ est inversible.
\begin{center}\texttt{-:-:-:-}\end{center}
% archive sheet number: 69

